\documentclass[../../main.tex]{subfiles}% 注意这里的文件路径不能用 ./main.tex ,否则用latexmk编译子文件会报错
\graphicspath{{\subfix{./image/}}} % 指定图片目录,后续可以直接使用图片文件名
% 注意这里的文件路径不能用 ../../image/ ,否则用latexmk编译子文件会报错

% 例如:
% \begin{figure}[H]
% \centering
% \includegraphics[scale=0.3]{图.png}
% \caption{}
% \label{figure:图}
% \end{figure}
% 注意:上述\label{}一定要放在\caption{}之后,否则引用图片序号会只会显示??.

\begin{document}

\section{点集的Lebesgue外测度}

\begin{definition}[Lebesgue外测度]
设 \(E \subset \mathbb{R}^n\). 若 \(\{I_k\}\) 是 \(\mathbb{R}^n\) 中的可数个开矩体, 且有
\begin{align*}
E \subset \bigcup_{k \geqslant  1} I_k,
\end{align*}
则称 \(\{I_k\}\) 为 \(E\) 的一个 \textbf{L-覆盖}(显然, 这样的覆盖有很多, 且每一个 \(L -\)覆盖 \(\{I_k\}\) 确定一个非负广义实值 \(\sum_{k \geqslant  1} |I_k|\) (可以是 \(+ \infty\), \(|I_k|\) 表示 \(I_k\) 的体积)). 称
\begin{align*}
m^*(E) = \inf \left\{ \sum_{k \geqslant  1} |I_k| : \{I_k\} \text{ 为 } E \text{ 的 } L-\text{覆盖} \right\}
\end{align*}
为点集 \(E\) 的 \textbf{Lebesgue 外测度}, 简称\textbf{外测度}.
\end{definition}
\begin{remark}
显然, 若 \(E\) 的任意的 \(L -\)覆盖 \(\{I_k\}\) 均有
\begin{align*}
\sum_{k \geqslant  1} |I_k| = + \infty,
\end{align*}
则 \(m^*(E) = + \infty\), 否则 \(m^*(E) < + \infty\). 
\end{remark}

\begin{definition}[集合上的外测度]
设 \(X\) 是一个非空集合, \(\mu^*\) 是定义在幂集 \(\mathcal{P}(X)\) 上的一个取广义实值的集合函数, 且满足：
\begin{enumerate}[(i)]
\item \(\mu^*(\varnothing)=0\), \(\mu^*(E) \geqslant  0\) (\(E \subset X\));
\item 若 \(E_1, E_2 \subset X\), \(E_1 \subset E_2\), 则 \(\mu^*(E_1) \leqslant  \mu^*(E_2)\);
\item 若 \(\{E_n\}\) 是 \(X\) 的子集列, 则有
\begin{align*}
\mu^*\left( \bigcup_{n = 1}^{\infty} E_n \right) \leqslant  \sum_{n = 1}^{\infty} \mu^*(E_n),
\end{align*}
\end{enumerate}
那么称 \(\mu^*\) 是 \(X\) 上的一个\textbf{外测度}.

若 \((X, d)\) 是一个距离空间, 且其上的外测度 \(\mu^*\) 还满足\textbf{距离外测度性质}:当 \(d(E_1, E_2)>0\) 时, 有
\begin{align*}
\mu^*(E_1 \cup E_2) = \mu^*(E_1) + \mu^*(E_2),
\end{align*}
那么称 \(\mu^*\) 是 \(X\) 上的一个\textbf{距离外测度}(利用距离外测度性质可以证明开集的可测性). 
\end{definition}

\begin{theorem}[\(\mathbb{R}^n\) 中点集的外测度性质]\label{theorem:R^n 中点集的外测度性质}
\begin{enumerate}[(1)]
\item 非负性: \(m^*(E) \geqslant  0\), \(m^*(\varnothing)=0\);

\item 单调性: 若 \(E_1 \subset E_2\), 则 \(m^*(E_1) \leqslant  m^*(E_2)\);

\item 次可加性: \(m^*\left(\bigcup_{k = 1}^{\infty} E_k\right) \leqslant  \sum_{k = 1}^{\infty} m^*(E_k)\).
\end{enumerate}
\end{theorem}
\begin{proof}
\begin{enumerate}[(1)]
\item 这可从定义直接得出.

\item 这是因为 \(E_2\) 的任一个 \(L -\)覆盖都是 \(E_1\) 的 \(L -\)覆盖.

\item 不妨设 \(\sum_{k = 1}^{\infty} m^*(E_k) < + \infty\). 对任意的 \(\varepsilon > 0\) 以及每个自然数 \(k\), 存在 \(E_k\) 的 \(L -\)覆盖 \(\{I_{k, l}\}\), 使得
\begin{align*}
E_k \subset \bigcup_{l = 1}^{\infty} I_{k, l}, \quad \sum_{l = 1}^{\infty} |I_{k, l}| < m^*(E_k) + \frac{\varepsilon}{2^k}.
\end{align*}
由此可知
\begin{align*}
\bigcup_{k = 1}^{\infty} E_k \subset \bigcup_{k, l = 1}^{\infty} I_{k, l}, \quad \sum_{k, l = 1}^{\infty} |I_{k, l}| \leqslant  \sum_{k = 1}^{\infty} m^*(E_k) + \varepsilon.
\end{align*}
显然, \(\{I_{k, l} : k, l = 1, 2, \cdots\}\) 是 \(\bigcup_{k = 1}^{\infty} E_k\) 的 \(L -\)覆盖, 从而有
\begin{align*}
m^*\left(\bigcup_{k = 1}^{\infty} E_k\right) \leqslant  \sum_{k = 1}^{\infty} m^*(E_k) + \varepsilon.
\end{align*}
由 \(\varepsilon\) 的任意性可知结论成立.
\end{enumerate}

\end{proof}

\begin{proposition}\label{proposition-R^n中的单点集的外测度为零}
\(\mathbb{R}^n\) 中的单点集的外测度为零, 即 \(m^*(\{x_0\}) = 0\), \(x_0 \in \mathbb{R}^n\).  同理, \(\mathbb{R}^n\) 中的点集
\[
\left\{ x = (\xi_1, \xi_2, \cdots, \xi_{i - 1}, t_0, \xi_i, \cdots, \xi_n) : a_j \leqslant  \xi_j \leqslant  b_j, j \neq i \right\}
\]
(\(n - 1\) 维超平面块)的外测度也为零.
\end{proposition}
\begin{proof}

这是因为可作一开矩体 \(I\), 使得 \(x_0 \in I\) 且 \(|I|\) 可任意地小.

\end{proof}

\begin{corollary}
若 \(E \subset \mathbb{R}^n\) 为可数点集, 则 \(m^*(E) = 0\).
\end{corollary}
\begin{remark}
由此可知有理点集的外测度 \(m^*(\mathbb{Q}^n)=0\). 这里我们看到了一个虽然处处稠密但外测度为零的可列点集.
\end{remark}
\begin{proof}

由外测度的次可加性不难证明.

\end{proof}

\begin{proposition}\label{proposition:Cantor集的外测度是零}
\([0,1]\) 中的 Cantor 集 \(C\) 的外测度是零.
\end{proposition}
\begin{remark}
这个\refpro{proposition:Cantor集的外测度是零}说明外测度为零的点集不一定是可列集.
\end{remark}
\begin{proof}

事实上, 因为 \(C = \bigcap_{n = 1}^{\infty} F_n\), 其中的 \(F_n\) (在构造 \(C\) 的过程中第 \(n\) 步所留存下来的) 是 \(2^n\) 个长度为 \(3^{-n}\) 的闭区间的并集, 所以我们有
\begin{align*}
m^*(C) \leqslant  m^*(F_n) \leqslant  2^n \cdot 3^{-n},
\end{align*}
从而得知 \(m^*(C)=0\). 

\end{proof}

\begin{proposition}
设 \(I\) 是 \(\mathbb{R}^n\) 中的开矩体,则  \(\overline{I}\) 是闭矩体,且\(m^*(I) = m^*(\overline{I}) = |I|\). 
\end{proposition}
\begin{proof}

显然$\overline{I}=I\cup \partial I$是闭矩体.先证$m^*(I)=|I|$.因为$I$是$I$自身的$L$-覆盖,所以$m^*(I)\leqslant |I|$.任取$I$的$L$-覆盖$\{I_k\}$,则
\[
\bigcup_{k = 1}^{\infty} I_{k} \supset I.
\]
从而由\nrefpro{proposition:矩体的性质}{(3)}可知
\begin{align*}
|I| \leqslant \sum_{k = 1}^{\infty} |I_k|,
\end{align*}
因此\(|I| \leqslant m^*(I)\).故$m^*(I)=|I|.$

再证$m^*(\overline{I})=|I|$.
对任给的 \(\varepsilon > 0\), 作一开矩体 \(J\), 使得 \(J \supset \overline{I}\) 且 \(|J| < |I| + \varepsilon\), 从而有
\begin{align*}
m^*(\overline{I}) \leqslant  |J| < |I| + \varepsilon.
\end{align*}
由 \(\varepsilon\) 的任意性可知 \(m^*(\overline{I}) \leqslant  |I|\). 现在设 \(\{I_k\}\) 是 \(\overline{I}\) 的任意的 \(L -\)覆盖, 则因为 \(\overline{I}\) 是有界闭集, 所以存在 \(\{I_k\}\) 的有限子覆盖
\[
\{I_{i_1}, I_{i_2}, \cdots, I_{i_l}\},\text{且}\bigcup_{j = 1}^l I_{i_j} \supset \overline{I}\supset I.
\]
于是由\nrefpro{proposition:矩体的性质}{(3)}可知
\begin{align*}
|I| \leqslant  \sum_{j = 1}^l |I_{i_j}| \leqslant  \sum_{k = 1}^{\infty} |I_k|,
\end{align*}
再由下确界是最大的下界可得 \(|I| \leqslant  m^*(\overline{I})\), 从而我们有 \(m^*(\overline{I}) = |I|\). 
综上,$m^*(I)=|I|=m^*(\overline{I})$.

\end{proof}

\begin{lemma}
设 \(E \subset \mathbb{R}^n\) 以及 \(\delta > 0\). 令
\begin{align*}
m^*_{\delta}(E) = \inf\left\{ \sum_{k = 1}^{\infty} |I_k| : \bigcup_{k = 1}^{\infty} I_k \supset E, \text{每个开矩体 } I_k \text{ 的边长} < \delta \right\},
\end{align*}
则 \(m^*_{\delta}(E)=m^*(E)\).
\end{lemma}
\begin{note}
这个引理告诉我们,今后可以对点集$E$的$L$-覆盖中的每个开矩体的边长做任意限制,而不影响$E$的外测度的值.
\end{note}
\begin{proof}

显然有 \(m^*_{\delta}(E) \geqslant  m^*(E)\). 为证明其反向不等式也成立，不妨设 \(m^*(E)< + \infty\). 由外测度的定义可知，对于任给的 \(\varepsilon > 0\)，存在 \(E\) 的 \(L -\)覆盖 \(\{I_k\}\)，使得
\begin{align*}
\sum_{k = 1}^{\infty} |I_k| \leqslant  m^*(E) + \varepsilon.
\end{align*}
对于每个 \(k\)，我们把 \(I_k\) 分割成 \(l(k)\) 个开矩体：
\[
I_{k,1}, I_{k,2}, \cdots, I_{k,l(k)},
\]
它们互不相交且每个开矩体的边长都小于 \(\frac{\delta}{2}\). 现在保持每个 \(I_{k,i}\) 的中心不动，边长扩大 \(\lambda(1 < \lambda < 2)\) 倍做出开矩体，并记为 \(\lambda I_{k,i}\)，显然，对每个 \(k\)，有
\begin{align*}
\bigcup_{i = 1}^{l(k)} \lambda I_{k,i} \supset I_k, \quad \sum_{i = 1}^{l(k)} |\lambda I_{k,i}| = \lambda^n \sum_{i = 1}^{l(k)} |I_{k,i}| = \lambda^n |I_k|.
\end{align*}
易知 \(\{\lambda I_{k,i} : i = 1,2,\cdots, l(k); k = 1,2,\cdots\}\) 是 \(E\) 的边长小于 \(\delta\) 的 \(L -\)覆盖，且有
\begin{align*}
\sum_{k = 1}^{\infty} \sum_{i = 1}^{l(k)} |\lambda I_{k,i}| = \lambda^n \sum_{k = 1}^{\infty} |I_k| \leqslant  \lambda^n (m^*(E) + \varepsilon),
\end{align*}
从而可知 \(m^*_{\delta}(E) \leqslant  \lambda^n (m^*(E) + \varepsilon)\). 令 \(\lambda \to 1\) 并注意到 \(\varepsilon\) 的任意性，我们得到 \(m^*_{\delta}(E) \leqslant  m^*(E)\). 这说明 \(m^*_{\delta}(E)=m^*(E)\).

\end{proof}

\begin{theorem}\label{theorem:距离大于零的两个点集的外测度满足可数可加性}
设 \(E_1, E_2\) 是 \(\mathbb{R}^n\) 中的两个点集. 若 \(d(E_1, E_2)>0\)，则
\begin{align*}
m^*(E_1 \cup E_2) = m^*(E_1) + m^*(E_2).
\end{align*}
\end{theorem}
\begin{proof}

由外测度的次可加性可知,只需证明 \(m^*(E_1 \cup E_2) \geqslant  m^*(E_1) + m^*(E_2)\) 即可. 为此，不妨设 \(m^*(E_1 \cup E_2)< + \infty\). 对任给的 \(\varepsilon > 0\)，作 \(E_1 \cup E_2\) 的 \(L -\)覆盖 \(\{I_k\}\)，使得
\begin{align*}
\sum_{k = 1}^{\infty} |I_k| < m^*(E_1 \cup E_2) + \varepsilon,
\end{align*}
其中 \(I_k\) 的边长都小于 \(\frac{d\left( E_1,E_2 \right)}{\sqrt{n}}\). 现在将 \(\{I_k\}\) 分为如下两组：
\begin{align}
(\text{i})J_{i_1}, J_{i_2}, \cdots,\bigcup_{k \geqslant  1} J_{i_k} \supset E_1 ; \quad (\text{ii})J_{l_1}, J_{l_2}, \cdots,\bigcup_{k \geqslant  1} J_{l_k} \supset E_2.\label{eq:::90jf43hg100.58}
\end{align}
且其中任一矩体皆不能同时含有 \(E_1\) 与 \(E_2\) 中的点.否则,不妨设存在$m\in [1,k]\cap \mathbb{N}$,使得$J_{i_m}$中同时含有$E_1$和$E_2$中的点.设$x_1\in J_{i_m}\cap E_1$,$x_2\in J_{i_m}\cap E_2$,则由$d(E_1,E_2)>0$可知
\begin{align*}
d(x_1,x_2) \geqslant d(E_1,E_2) > 0.
\end{align*}
又因为\(I_k\) 的边长都小于 \(\frac{d\left( E_1,E_2 \right)}{\sqrt{n}}\),所以由\nrefpro{proposition:矩体的性质}{(1)}可知
\begin{align*}
d(x_1,x_2)\leqslant \mathrm{diam}\left( J_{i_m} \right) \leqslant \sqrt{\left( \sum_{i=1}^n{\frac{d\left( E_1,E_2 \right)}{\sqrt{n}}} \right) ^2}=d(E_1,E_2)<d(x_1,x_2).
\end{align*}
上式显然矛盾!
故\eqref{eq:::90jf43hg100.58}式成立.
从而得
\begin{align*}
m^*(E_1 \cup E_2) + \varepsilon > \sum_{k \geqslant  1} |I_k| = \sum_{k \geqslant  1} |J_{i_k}| + \sum_{k \geqslant  1} |J_{l_k}| \geqslant  m^*(E_1) + m^*(E_2).
\end{align*}
再由 \(\varepsilon\) 的任意性可知 \(m^*(E_1 \cup E_2) \geqslant  m^*(E_1) + m^*(E_2)\). 

\end{proof}

\begin{corollary}\label{corollary:距离大于零的两个点集的外测度满足可数可加性}
设 \(E_1, E_2,\cdots,E_n\) 是 \(\mathbb{R}^n\) 中的$n$个点集. 若 \(d(E_i, E_j)>0\)($i\ne j$)，则
\begin{align*}
m^*\left( \bigcup_{i=1}^n{E_i} \right) =\sum_{i=1}^n{m^*\left( E_i \right)}.
\end{align*}
\end{corollary}
\begin{proof}

当$n=1$时结论显然成立.假设当$n=k$时结论成立,现在考虑$n=k+1$的情形.由\hyperref[proposition:点集间的距离的性质]{点集间的距离的性质}及$d(E_i,E_j)>0$（$i\neq j$）可知
\begin{align*}
d\left(E_{k + 1},\bigcup_{i = 1}^kE_i\right)=\min_{i = 1,2,\cdots,k}d\left(E_{k + 1},E_i\right)>0.
\end{align*}
故再由\refthe{theorem:距离大于零的两个点集的外测度满足可数可加性}和归纳假设可得
\begin{align*}
m^*\left(\bigcup_{i = 1}^{k + 1}E_i\right)&=m^*\left(E_{k + 1}\cup\bigcup_{i = 1}^kE_i\right)=m^*\left(E_{k + 1}\right)+m^*\left(\bigcup_{i = 1}^kE_i\right)\\
&=m^*\left(E_{k + 1}\right)+\sum_{i = 1}^km^*(E_i)
=\sum_{i = 1}^{k + 1}m^*(E_i).
\end{align*}
因此由数学归纳法可知结论成立。 

\end{proof}

\begin{theorem}[外测度的平移不变性]\label{theorem:外测度的平移不变性}
设 \(E \subset \mathbb{R}^n\), \(x_0 \in \mathbb{R}^n\). 记 \(E + \{x_0\} = \{x + x_0, x \in E\}\), 则
\begin{align}
m^*(E + \{x_0\}) = m^*(E). \label{eq:2.1}
\end{align}
\end{theorem}
\begin{remark}
对集合做相同的平移并不会改变集合之间的关系(交、并、差、补、子集等).
\end{remark}
\begin{proof}

首先, 对于 \(\mathbb{R}^n\) 中的开矩体 \(I\), 易知 \(I + \{x_0\}\) 仍是一个开矩体且其相应边长均相等, \(|I| = |I + \{x_0\}|\). 其次, 对 \(E\) 的任意的 \(L -\)覆盖 \(\{I_k\}\), \(\{I_k + \{x_0\}\}\) 仍是 \(E + \{x_0\}\) 的 \(L -\)覆盖. 从而由
\begin{align*}
m^*(E + \{x_0\}) \leqslant  \sum_{k = 1}^{\infty} |I_k + \{x_0\}| = \sum_{k = 1}^{\infty} |I_k|
\end{align*}
可知(对一切 \(L -\)覆盖取下确界)
\begin{align*}
m^*(E + \{x_0\}) \leqslant  m^*(E).
\end{align*}
反之, 考虑对 \(E + \{x_0\}\) 作向量 \(-x_0\) 的平移, 可得原点集 \(E\). 同理又有
\begin{align*}
m^*(E) \leqslant  m^*(E + \{x_0\}).
\end{align*} 

\end{proof}

\begin{theorem}[外测度的数乘]\label{theorem:外测度的数乘}
设 \(E \subset \mathbb{R}\), \(\lambda \in \mathbb{R}\), 记 \(\lambda E = \{ \lambda x : x \in E\}\), 则
\begin{align*}
m^*(\lambda E) = |\lambda| m^*(E).
\end{align*}
\end{theorem}
\begin{proof}

因为 \(E \subset \bigcup_{n \geqslant  1} (a_n, b_n)\) 等价于 \(\lambda E \subset \bigcup_{n \geqslant  1} \lambda (a_n, b_n)\), \(m^*([a_n, b_n]) = m^*((a_n, b_n))\), 且对任一区间 \((\alpha, \beta)\), 有
\begin{align*}
m^*(\lambda (\alpha, \beta)) = |\lambda| m^*((\alpha, \beta)) = |\lambda| (\beta - \alpha),
\end{align*}
所以按外测度定义可得 \(m^*(\lambda E) = |\lambda| m^*(E)\). 

\end{proof}

\begin{proposition}
\begin{enumerate}[(1)]
\item 设 $A\subseteq\mathbf{R}^n$ 且 $m^*(A)=0$,则对任意的 $B\subseteq\mathbf{R}^n$,有
\begin{align*}
m^*(A\cup B)=m^*(B)=m^*(B\setminus A).
\end{align*}
\item 设 $A,B\subseteq\mathbf{R}^n$,且 $m^*(A),m^*(B)<\infty$,则
\begin{align*}
|m^*(A)-m^*(B)|\leqslant m^*(A\triangle B);
\end{align*}
\item 设 $A,B\subseteq\mathbf{R}^n$.若 $m^*(A\triangle B)=0$,则 $m^*(A)=m^*(B)$.
\item 设 $A,B$ 与 $C$ 是 $\mathbf{R}^n$ 中的点集,且有 $m^*(A\triangle B)=0$, $m^*(B\triangle C)=0$,则 $m^*(A\triangle C)=0$.
\item 设 $E\subseteq\mathbf{R}^n$.若对任意的 $x\in E$,存在开球 $B(x,\delta_x)$,使得 $m^*(E\cap B(x,\delta_x))=0$,则 $m^*(E)=0$.
\end{enumerate}
\end{proposition}
\begin{proof}
\begin{enumerate}[(1)]
\item 注意,我们有不等式:
\begin{align*}
m^*(A\cup B)\leqslant m^*(A)+m^*(B)=m^*(B)\leqslant m^*(A\cup B),
\end{align*}
\begin{align*}
m^*(B)\leqslant m^*(B\setminus A)+m^*(B\cap A)\leqslant m^*(B\setminus A)+m^*(A)=m^*(B\setminus A)\leqslant m^*(B).
\end{align*}

\item 因为 $A\subseteq B\cup(A\triangle B)$,所以 $m^*(A)\leqslant m^*(B)+m^*(A\triangle B)$.从而可知 $m^*(A)-m^*(B)\leqslant m^*(A\triangle B)$.类似地,又可得 $m^*(B)-m^*(A)\leqslant m^*(A\triangle B)$.综合此两结论,即得所证.

\item 由 $m^*(A)\leqslant m^*(B)+m^*(A\triangle B)=m^*(B)$ 可知,$m^*(A)\leqslant m^*(B)$.又由 $m^*(B)\leqslant m^*(A)+m^*(A\triangle B)$ 可得 $m^*(B)\leqslant m^*(A)$.

\item 注意到公式
\begin{align*}
(A\cup B\cup C)\setminus(A\cap B\cap C)=(A\triangle B)\cup(B\triangle C),\\
A\cap C\supseteq A\cap B\cap C,\quad A\cup C\subseteq A\cup B\cup C,
\end{align*}
可知 $(A\cup C)\setminus(A\cap C)\subseteq(A\cup B\cup C)\setminus(A\cap B\cap C)$.从而有
\begin{align*}
m^*(A\triangle C)\leqslant m^*(A\triangle B)+m^*(B\triangle C)=0.
\end{align*}
\item 依题设可知$E\subseteq \bigcup_{x\in E}{B\left( x,\delta _x \right)}$,由\hyperref[theorem:Lindelöf定理]{Lindelöf定理}知,存在 $E$ 的可数覆盖球列 $\{B_k\triangleq B(x_k,\delta_{x_k})\}$,使得 $E\subseteq\bigcup\limits_{k=1}^{\infty}B_k$,且 $m^*(E\cap B_k)=0$.从而知
\begin{align*}
E=\bigcup\limits_{k=1}^{\infty}(E\cap B_k),\quad m^*(E)\leqslant\sum\limits_{k=1}^{\infty}m^*(E\cap B_k)=0.
\end{align*}
\end{enumerate}

\end{proof}

\begin{proposition}
\begin{enumerate}[(1)]
\item 设 $E\subseteq\mathbf{R}^n$,试证明存在 $G_\delta$ 集 $\widetilde{G}$: $\widetilde{G}\supseteq E$ 且 $m^*(\widetilde{G})=m^*(E)$.

\item 设 $E_k\subseteq\mathbf{R}^n$ ($k\in\mathbf{N}$).若 $\sum\limits_{k=1}^{\infty}m^*(E_k)<+\infty$,则 $m\left(\varlimsup_{k\to\infty}E_k\right)=0$.
\end{enumerate}
\end{proposition}
\begin{proof}
\begin{enumerate}[(1)]
\item 依定义,对 $k\in\mathbf{N}$,存在矩体列 $\{I_{k,i}\}$,使得 $\bigcup\limits_{i=1}^{\infty}I_{k,i}\supseteq E$,且有
\begin{align*}
\sum\limits_{i=1}^{\infty}|I_{k,i}|\leqslant m^*(E)+\frac{1}{k}\quad (k=1,2,\cdots).
\end{align*}
令 $\widetilde{G}_k=\bigcup\limits_{i=1}^{\infty}I_{k,i}$ ($k\in\mathbf{N}$), $\widetilde{G}=\bigcap\limits_{k=1}^{\infty}\widetilde{G}_k$,易知 $\widetilde{G}\supseteq E$ 且 $\widetilde{G}$ 是 $G_\delta$ 集.我们有(对 $k\in\mathbf{N}$)
\begin{align*}
m^*(E)\leqslant m^*(\widetilde{G})\leqslant m^*(\widetilde{G}_k)\leqslant \sum\limits_{i=1}^{\infty}|I_{k,i}|\leqslant m^*(E)+\frac{1}{k},
\end{align*}
令 $k\to\infty$ 即得 $m^*(E)=m^*(\widetilde{G})$.
\item 注意 $\varlimsup_{k\to\infty}E_k=\bigcap\limits_{m=1}^{\infty}\bigcup\limits_{k=m}^{\infty}E_k$,且依题设知,对任给 $\varepsilon>0$,存在 $N$,使得 $\sum\limits_{k=N}^{\infty}m^*(E_k)<\varepsilon$.从而对任意 $j\in\mathbf{N}$,有
\begin{align*}
m^*\left(\varlimsup_{k\to\infty}E_k\right)\leqslant m^*\left(\bigcup\limits_{k=j}^{\infty}E_k\right)\leqslant \sum\limits_{k=j}^{\infty}m^*(E_k).
\end{align*}
由此知 $m^*\left(\varlimsup_{k\to\infty}E_k\right)\leqslant \sum\limits_{k=N}^{\infty}m^*(E_k)<\varepsilon$.
\end{enumerate}

\end{proof}

\begin{definition}[内测度]
设$\mathscr{R}$为$X$的某些子集所成的环，$\mu$为$\mathscr{R}$上的测度。则$\mathscr{K}(\mathscr{R})$上的集函数
\[
\mu_*(E) = \sup\{\mu(F) \mid E \supseteq F \in \mathscr{R}\}
\]
称$\mu_*$为\textbf{内测度}.
\end{definition}

\begin{theorem}[内测度的性质]\label{theorem:内测度的基本性质}
$\mu_*$为\textbf{内测度}具有下列各性质：
\begin{enumerate}[(1)]
\item 非负性：$\mu_*(E)\geqslant 0$。特别地，$\mu_*(\varnothing)=0$；
\item 单调性：若$E_1,E_2\in\mathscr{K}(\mathscr{R})$，$E_1\subseteq E_2$，则$\mu_*(E_1)\leqslant \mu_*(E_2)$；
\item 若$E\in\mathscr{R}$，则$\mu_*(E)=\mu(E)$；
\item 若$E_n\in\mathscr{K}(\mathscr{R})$，$E_i\cap E_j=\varnothing$，$i\neq j$，则
\begin{align*}
\mu_*\left(\bigcup\limits_{n=1}^{m}E_n\right)\geqslant \sum\limits_{n=1}^{m}\mu_*(E_n).
\end{align*}
进而，有
\begin{align*}
\mu_*\left(\bigcup\limits_{n=1}^{\infty}E_n\right)\geqslant \sum\limits_{n=1}^{\infty}\mu_*(E_n);
\end{align*}
\item 对$E\in\mathscr{R}$，$F\in\mathscr{K}(\mathscr{R})$，有
\begin{align*}
\mu_*(F) = \mu_*(F\cap E) + \mu_*(F\setminus E) 
= \mu_*(F\cap E) + \mu_*(F\cap E^c);
\end{align*}
\item 若$E\in\mathscr{K}(\mathscr{R})$，则
\begin{align*}
\mu_*(E)\leqslant \mu^*(E).
\end{align*}
\end{enumerate}
\end{theorem}
\begin{proof}

根据$\mu_*$的定义和$\mu$的性质，显然有(1)、(2)、(3)。

(4) 如果有$\mu_*(E_n)=+\infty$，则$\mu_*\left(\bigcup\limits_{n=1}^{m}E_n\right)\overset{(2)}{\geqslant} \mu_*(E_n)=+\infty$，从而
\begin{align*}
\mu_*\left(\bigcup\limits_{n=1}^{m}E_n\right)=+\infty = \sum\limits_{n=1}^{m}\mu_*(E_n).
\end{align*}
如果$\mu_*(E_n)<+\infty$，$n=1,2,\cdots,m$，则对$E_i,E_j\in\mathscr{K}(\mathscr{R})$，$E_i\cap E_j=\varnothing$，$i\neq j$和$\forall\varepsilon>0$，取$F_n$，使得$E_n\supseteq F_n\in\mathscr{R}$及
\begin{align*}
\mu_*(E_n) < \mu(F_n) + \frac{\varepsilon}{2^n}.
\end{align*}
于是，$\bigcup\limits_{n=1}^{m}E_n \supseteq \bigcup\limits_{n=1}^{m}F_n$。再从(2) 立即得到
\begin{align*}
\mu_*\left(\bigcup\limits_{n=1}^{m}E_n\right) &\geqslant \mu_*\left(\bigcup\limits_{n=1}^{m}F_n\right) = \mu\left(\bigcup\limits_{n=1}^{m}F_n\right) = \sum\limits_{n=1}^{m}\mu(F_n) \\
&> \sum\limits_{n=1}^{m}\mu_*(E_n) - \sum\limits_{n=1}^{m}\frac{\varepsilon}{2^n} > \sum\limits_{n=1}^{m}\mu_*(E_n) - \varepsilon.
\end{align*}
令$\varepsilon\to 0^+$，有
\begin{align*}
\mu_*\left(\bigcup\limits_{n=1}^{m}E_n\right)\geqslant \sum\limits_{n=1}^{m}\mu_*(E_n).
\end{align*}
进而，如果$E_n\in\mathscr{K}(\mathscr{R})$，$E_i\cap E_j=\varnothing$，$i\neq j$，$n,i,j=1,2,\cdots$，则由上结论，知
\begin{align*}
\mu_*\left(\bigcup\limits_{n=1}^{\infty}E_n\right)\geqslant \mu_*\left(\bigcup\limits_{n=1}^{m}E_n\right)\geqslant \sum\limits_{n=1}^{m}\mu_*(E_n).
\end{align*}
令$m\to+\infty$得到
\begin{align*}
\mu_*\left(\bigcup\limits_{n=1}^{\infty}E_n\right)\geqslant \sum\limits_{n=1}^{\infty}\mu_*(E_n).
\end{align*}

(5) 对$E\in\mathscr{R}$，$F\in\mathscr{K}(\mathscr{R})$，由于
\[
F = (F\cap E)\cup (F\setminus E)
\]
且
\[
(F\cap E)\cap (F\setminus E) = \varnothing
\]
和(4)有
\begin{align*}
\mu_*(F)\geqslant \mu_*(F\cap E) + \mu_*(F\setminus E).
\end{align*}
再证$\mu_*(F)\leqslant \mu_*(F\cap E) + \mu_*(F\setminus E)$。如果$\mu_*(F\cap E)=+\infty$或$\mu_*(F\setminus E)=+\infty$，则必有$\mu_*(F)\geqslant \mu_*(F\cap E) + \mu_*(F\setminus E)=+\infty$，$\mu_*(F)=+\infty=\mu_*(F\cap E) + \mu_*(F\setminus E)$。如果$\mu_*(F\cap E)<+\infty$，$\mu_*(F\setminus E)<+\infty$，则对$\forall\varepsilon>0$，由$\mu_*$定义，$\exists F_0\in\mathscr{R}$使得$F_0\subseteq F$且$\mu_*(F)<\mu(F_0)+\varepsilon$。又$\mu(F_0)=\mu(F_0\cap E)+\mu(F_0\setminus E)$，故
\begin{align*}
\mu_*(F) &< \mu(F_0) + \varepsilon = \mu(F_0\cap E) + \mu(F_0\setminus E) + \varepsilon \\
&\leqslant \mu_*(F\cap E) + \mu_*(F\setminus E) + \varepsilon,
\end{align*}
令$\varepsilon\to 0^+$得到
\begin{align*}
\mu_*(F)\leqslant \mu_*(F\cap E) + \mu_*(F\setminus E).
\end{align*}
综合上述，有
\begin{align*}
\mu_*(F) = \mu_*(F\cap E) + \mu_*(F\setminus E).
\end{align*}

(6) 设$E\in\mathscr{K}(\mathscr{R})$，对$\forall F\in\mathscr{R}$，$E_i\in\mathscr{R}$，$i\in\mathbb{N}$，如果$F\subseteq E\subseteq \bigcup\limits_{i=1}^{\infty}E_i$，则
\begin{align*}
\mu(F) = \mu^*(F)\leqslant \mu^*\left(\bigcup\limits_{i=1}^{\infty}E_i\right)\leqslant \sum\limits_{i=1}^{\infty}\mu^*(E_i) = \sum\limits_{i=1}^{\infty}\mu(E_i).
\end{align*}
由此推得
\begin{align*}
\mu_*(E) = \inf\{\mu(F) \mid E\supseteq F\in\mathscr{R}\} \leqslant \sum\limits_{i=1}^{\infty}\mu(E_i), \\
\mu_*(E)\leqslant \inf\left\{\sum\limits_{i=1}^{\infty}\mu(E_i) \mid E_i\in\mathscr{R}, E\subseteq \bigcup\limits_{i=1}^{\infty}E_i\right\} = \mu^*(E).
\end{align*}

\end{proof}








\end{document}